%% Supplementary material for
%%   Prediction-interval coverage is a design-based quantity, and clustered
%%   sampling breaks it
%% Environmental Technology & Innovation (Elsevier)
%% Build: pdflatex -> bibtex -> pdflatex -> pdflatex
\documentclass[preprint,12pt,authoryear]{elsarticle}

\usepackage{amsmath,amssymb}
\usepackage{graphicx}
\usepackage{booktabs}
\usepackage{multirow}
\usepackage[T1]{fontenc}
\usepackage[utf8]{inputenc}
\usepackage[hidelinks,bookmarks=false,pdfpagemode=UseNone]{hyperref}
\hypersetup{colorlinks=false,pdfborder={0 0 0}}
\usepackage{microtype}
\usepackage{url}
\urlstyle{same}
\def\UrlBreaks{\do\_\do\/\do\.\do\-}
\newcommand{\file}[1]{\path{#1}}
\setlength{\emergencystretch}{3em}
\setlength{\tabcolsep}{4.5pt}

\newcommand{\Mgha}{Mg\,ha\textsuperscript{-1}}

%% Everything here is numbered S1, S2, ... and Appendix A, so nothing collides
%% with the main text.
\renewcommand{\thetable}{S\arabic{table}}
\renewcommand{\thefigure}{S\arabic{figure}}
\renewcommand{\theequation}{A.\arabic{equation}}
\renewcommand{\thesection}{A}

\begin{document}

\begin{center}
{\large\bfseries Supplementary material}\\[2pt]
{\itshape Prediction-interval coverage is a design-based quantity, and clustered
sampling breaks it}\\[2pt]
Hugo Rodrigues, Golmar Golmohammadi
\end{center}

\vspace{1em}

\section*{Supplementary tables}

\begin{table}[!ht]
\centering
\caption{Comparison against the untuned default forest under four estimands,
simple random calibration design. RMSE in \Mgha{}.}
\label{stab:estimands}
\small
\begin{tabular}{lrrrrr}
\toprule
Estimand & Reps & RF & This method & $\Delta$RMSE & Wins \\
\midrule
\multicolumn{6}{l}{\emph{Design-based map accuracy}} \\
Population (SRS 20{,}000) & 10 & 32.18 & 31.51 & $-0.669$ & 10/10 \\
Design-based (SRS 500)    & 30 & 32.49 & 31.62 & $-0.869$ & 28/30 \\
\midrule
\multicolumn{6}{l}{\emph{Model comparison only, not map accuracy}} \\
Random 10-fold CV         & 10 & 32.24 & 31.12 & $-1.125$ & 10/10 \\
Spatial block CV (350 km) & 10 & 34.72 & 33.45 & $-1.275$ & 10/10 \\
\bottomrule
\end{tabular}
\end{table}

\begin{table}[!ht]
\centering
\caption{Selected corrector across all 90 repetitions.}
\label{stab:modes}
\small
\begin{tabular}{lr}
\toprule
Selected corrector & Count \\
\midrule
Residual random forest & 72 \\
Ensemble of residual learners & 10 \\
Ordinary kriging of residuals & 3 \\
Cubist delta & 1 \\
Linear recalibration & 1 \\
Residual Cubist & 1 \\
Single-hidden-layer network & 1 \\
Kernel (radial basis) smoother & 1 \\
\bottomrule
\end{tabular}

\end{table}

\begin{table}[!ht]
\centering
\caption{Simulation S1. Proportion of runs in which anisotropy is declared, by
true ratio and signal-to-noise ratio, under simple random sampling with 500
points. A true ratio of 1 is an isotropic field.}
\label{stab:sim1}
\small
\begin{tabular}{rrr}
\toprule
True ratio $\varrho$ & Signal-to-noise & Anisotropy declared \\
\midrule
1.00 & 2 & 0.00 \\
0.60 & 2 & 0.00 \\
0.35 & 2 & 0.29 \\
0.20 & 2 & 0.92 \\
1.00 & 8 & 0.00 \\
0.60 & 8 & 0.00 \\
0.35 & 8 & 0.58 \\
0.20 & 8 & 1.00 \\
\bottomrule
\end{tabular}

\end{table}

\begin{table}[!ht]
\centering
\caption{Simulation S2. Proportion of runs declaring anisotropy when the true
geometry is held fixed and only the sampling design changes.}
\label{stab:sim2}
\small
\begin{tabular}{lrr}
\toprule
Sampling design & True ratio $\varrho$ & Anisotropy declared \\
\midrule
Clustered & 0.35 & 0.500 \\
Simple random & 0.35 & 0.583 \\
Clustered & 1.00 & 0.083 \\
Simple random & 1.00 & 0.000 \\
\bottomrule
\end{tabular}

\end{table}

\clearpage

\section{Full specification of the procedure}
\label{app:spec}

Every setting needed to reproduce the experiments is stated here; all values are
the defaults in the released code unless noted.

\paragraph{Cross-fitting and selection} $K=5$ folds assigned by
\texttt{caret::createFolds} on the response. The shrinkage grid is
$\lambda \in [-0.25, 0.60]$ in steps of $0.05$, 18 values, for $n \ge 300$;
narrower grids apply at smaller $n$. Selection minimises Eq.~(8) over the
correction library crossed with that grid, with $\lambda=0$ included.

\paragraph{Base forest profiles} Six fixed profiles are compared on identical
folds: \file{default_rf}, \file{more_trees}, \file{deeper_trees},
\file{subsample}, \file{wide_mtry}, \file{wide_deep_subsample}. Across
them the number of trees takes values in $\{500,1000,1500,2000\}$,
\texttt{mtry} in $\{5,9,14\}$ for $p=28$, minimum node size in $\{1,5\}$ and
sample fraction in $\{0.632,0.8,1\}$; residual forests use 300 or 500 trees.
Where explicit grid tuning is enabled instead, the base grid is capped at 18
candidates and the residual grid at 12, scored out of bag.

\paragraph{Correction library} Residual random forest; residual Cubist with 10
committees and 0 neighbours; a single-hidden-layer network; ordinary kriging of
residuals; a radial basis or kernel smoother; the unweighted mean of the
residual learners; and bias, linear, quadratic and stacked recalibrations. The
radial basis places centres by $k$-means on the coordinates, their number
derived from the domain area, the practical range and the nugget ratio and
clamped to $[8,48]$, with bandwidth $b=m_b\rho$.

\paragraph{Variogram} Empirical variogram with the robust estimator of
\citet{cressie1980}, cutoff at half the maximum pairwise distance, 15 lags.
Exponential, spherical, Gaussian and Mat\'ern models are each fitted by weighted
least squares to the same empirical variogram, and the simplest family within a
relative tolerance of $0.05$ of the best weighted residual sum of squares is
retained. The practical range is the distance at which the fitted model first
reaches 95 per cent of its sill, evaluated numerically so that it means the same
thing for every family.

\paragraph{Search axes} Three variogram configurations join the default: cutoff
$0.35$ with 15 lags, cutoff $0.65$ with 20 lags, and the classical
\citet{matheron1963} estimator at the default binning. The range multiplier
takes $m_b \in \{0.25,0.5,1.0\}$, with $0.5$ the default arm. The anisotropic
variant is offered for the two leading profiles; candidate geometries are angles
$\{0,22.5,\dots,157.5\}$ degrees crossed with ratios
$\{1,0.75,0.55,0.40,0.25\}$, scored by 5-fold out-of-fold error of a kernel
smoother whose bandwidth is $0.15$ times the median pairwise distance in the
transformed space, and accepted only on an improvement exceeding 3 per cent.

\paragraph{Gate and intervals} The gate length multiplier is chosen by spatially
blocked cross-validation with four $k$-means blocks, the candidate set including
the value that disables the gate. Intervals use CV+ with $K=5$ and
$\alpha=0.10$, on a calibration fold split drawn independently of the selection
split. The covariate-shift variant estimates density ratios with a forest
classifier separating calibration points from a sample of the population
covariates, clipping weights at 20; the spatial-block variant refits on five
$k$-means blocks. The quantile regression forest baseline uses \texttt{ranger}
with \texttt{quantreg = TRUE} and the same predictors.

\paragraph{Sampling designs} Simple random sampling draws $n=500$ cells without
replacement; the systematic design draws a regular grid; the clustered design is
two-stage with 20 primary units chosen over 100 $k$-means centres and 25 points
each. The validation sample is an independent simple random sample of $n=500$;
the Population estimand uses 20{,}000 cells drawn after excluding the
calibration cells. Spatial block cross-validation uses complete-linkage
clustering of coordinates cut at 350 km.

\paragraph{Preprocessing, missing data, seeds} Covariates and response enter
untransformed. Rows with non-finite response or coordinates are dropped before
variogram fitting, and constant or duplicated predictor columns are removed
before fitting. Every \texttt{ranger} call receives an explicit seed and
\texttt{ranger} 0.18 and later is reproducible across thread counts; the
calibration draw for repetition $i$ uses seed $123 + 1000(i-1)$. The internal
random state is restored on exit with \texttt{withr}, so calling the estimator
does not alter the user's stream.

\paragraph{Software and scripts} R 4.6.0 with \texttt{ranger}
\citep{wright2017}, \texttt{gstat} \citep{pebesma2004}, \texttt{caret},
\texttt{Cubist} and \texttt{nnet}. The estimator is the \texttt{GeoVersa}
package. Every table and figure in this paper is produced by
\file{run_dbtarf_wadoux_comparison.R} (accuracy campaigns),
\file{run_dbtarf_protocols_comparison.R} (estimands),
\file{simulations.R} (the simulation section of the main text), \file{tables_manuscript.R}
(all tables) and \file{maps_compute.R} with \file{maps_plot.R} (the maps).



\bibliographystyle{elsarticle-harv}
\bibliography{references}

\end{document}
